\chapter{Predicate Constructions}\label{ch:predicates}

\section{Stative verbs}\label{sec:stative}

Stative predicates form a morphosyntactic subclass within the verbal system.
They describe a condition or property holding of a participant, require ordinary
finite Asp.Align, and need no copula. Their lower \synhead{StateP} is distinct
from the eventive \synhead{BecomeP} introduced by overt inchoative morphology.
This distinction does not require a universal adjective category. The comparative evidence motivates
the separation of state and transition, rather than Iridian's particular
exponents: Japanese and Korean resist rate-adverb-triggered inchoative
readings of simple property states and instead employ overt change-of-state
predicates
\cite{smithhopperdietzelkoontzgarboden2025}. Korean categorial arguments
remain compatible with recognising language-specific descriptive classes
without declaring every property predicate universally adjectival
\cite{kimaeryung2024}.

\subsection{State types and aspectual interpretation}\label{sec:stative-stability}

Four classes distinguish lower predicate meanings without imposing a new
universal adjective category. Quality statives include \ird{mužn-},
\ird{hradn-} and \ird{torban-}; dimensional statives include
\ird{svorn-} and \ird{četn-}; experiential and condition statives include
\ird{ptakn-}, \ird{smrn-} and \ird{všihn-}; relational statives include
\ird{tvírzn-} and \ird{barocezn-}, derived through \ird{-(e)zn-}.
Lexical positional, possessive and locative states remain available outside
productive \ird{-n-} formation. These classes condition scale structure,
participant interpretation and contextual aspectual licensing; they are not
four automatic inflectional paradigms.


Lexical specification keeps these distinctions independent. An entry records
its selected stem and quantitative class, eligible nucleus and protected
quantity, argument frame, gradability, licensed aspectual construals,
Become class and availability, Manifest availability, ordinary attributive
class, manner formation and participant nominalisation. Psychological entries
add separate stimulus and perspective flags. A property may be gradable
without admitting a causal transition, and a short stem may lack either
derivational head. Class membership is therefore not inferred from the four
semantic labels or from vowel length.

Three argument-frame types must be distinguished: an experiencer alone;
a lexically selected experiencer with a stimulus or content relation; and an
experiencer whose stimulus is licensed by overt Appl. Distinguishing
experiencer-only and stimulus-bearing frames has a comparison in Korean
psych-predicate analysis \cite{kimyoonshin2026}. The listed
\ird{všihn-} paradigm has an experiencer-only base and an optional stimulus
construction in \ird{-éb-}, not a freely selected additional finite object.
The construction supplies content about which the experience holds. A stimulus
is not thereby an actual cause, and an argument-frame label never supplies
an unexpressed Caus head. No additional bare two-place psychological frame is
assigned merely to fill this three-way taxonomy.

Characteristic qualities favour the simple continuous. A temporary condition
also remains a state: its temporal limitation alone does not assert entry
into it. Marked expressive or stage-level contexts can license a simple-state
progressive; a bounded occasion can license a simple-state perfective.
Retrospective presents an anterior state without entailing that its beginning
was witnessed, caused or described, or that it subsequently ceased. A projected
contemplative state may already obtain; projection is not an assertion of
non-inception. Overt Become supplies the ordinary
transition reading in each class where change into the specified state is
semantically coherent. Relational membership need not have a natural open
scale; its transition requires a context that fixes the relevant relation.

Perspective-sensitive entries distinguish an internal experiencer from the
perspective holder licensing an unmediated experiential assertion. The default
anchor in an independent declarative is the speaker. Direct continuous
\ird{všihna} can then describe an internally accessible experience.
Questions can establish an addressee's experiential perspective; reported
thought, a licensed attitude complement and narrative perspective can establish
another anchor. A question or an embedded clause does not automatically shift
that anchor: its clause type, attitude relation or discourse must independently
identify the relevant perspective holder.

The dependency is interpreted with the clause's existing Mood, Evid and attitude
structure, rather than realised by an additional universal perspective affix.
The lower lexical predicate supplies its perspective-sensitive interpretation;
Mood and clause type constrain the speech-act participant, and a licensed
attitude or evidential interpretation can mediate access to that experience.
A Mood-based account of Korean perspective effects supplies a comparison
\cite{park2026perspective}, without establishing an identical Iridian
analysis. This does not merge an unrestricted new projection between existing heads.
Comparative analyses of Korean experiential and perspective-sensitive
constructions support separating perspective from grammatical person
\cite{nam2024,sakaba2026}; they do not determine the Iridian head order.

Indirect Evid marks the source of knowledge for the proposition, not a
psychological stem class. Thus \ird{všihna-l} gives \ird{všihnal},
linked \ird{všihnalě}, for a suitably reported or inferred state. A third-person
experiencer does not automatically require \ird{-al}/\ird{-l}, and a
reported or shifted experiential perspective does not itself assert visible
behaviour. Existing evidential allomorphy and its position after aspect and
Mood remain unchanged (§\,\ref{sec:vm-evid}). Internally felt bodily
conditions require an individually specified perspective-sensitive reading;
an externally observable condition does not inherit the restriction. Beauty,
size and colour do not require evidential morphology merely because they are
states. Outward display is separately derived by the restricted Manifest
construction (§\,\ref{sec:stative-manifest}); it does not replace the source
or perspective dependency of an assertion about genuine feeling.

\subsection{Valency operations with stative verbs}\label{sec:stative-valency}

Stative causatives describe the bringing about of a state or a change along
the scale expressed by the base. The progressive presents that change as
unfolding, whereas the perfective presents a bounded causing event; neither
viewpoint specifies whether the causer acts directly or through an intermediary.
Stative stems in \ird{-n-} select expanded causative \ird{-er-}, as in
\ird{svorn-ér-ime}. Overt outer Voice independently takes \ird{-an-}
after the causative, giving perfective \ird{svorn-ér-an-ik} and
retrospective \ird{svorn-ér-an-av}. The stem-final stative \ird{-n-},
the rhotic causative, and outer orientation therefore occupy distinct
positions, without metathesis or a joint causative--Voice exponent.

\pex\label{ex:val-stative-caus}
\a
\begingl
\gla Marek problemam svornérime.//
\glb Marek.\Dir{} problem.\Ind{} be.big-\Caus{}-\Prog{}//
\glft \trsl{Marek is making the problem bigger.}//
\endgl
\a
\begingl
\gla Problem Marku svornéranik.//
\glb problem.\Dir{} Marek.\Trs{} be.big-\Caus{}-\textsc{trs}-\Pf{}//
\glft \trsl{Marek made the problem bigger.}//
\endgl
\xe

Causative can also select an already inchoative predicate. Default compact
\ird{-r-} follows liquid-final Become without changing its protected quantity:
\ird{mužnél-r-ime} \trsl{is causing to become beautiful} and
\ird{mužnél-r-an-ik} \trsl{caused to become beautiful}, with the usual
NOM/ERG causee frames. The sequence /l.r/ permits a coda--onset division;
coda /l/ vocalises by the ordinary rule, without deletion or an invented
cluster repair. Protected \ird{tvírznel-r-ime} behaves likewise, where
change into the relevant relation is semantically licensed. Become introduces
the transition and Caus the causer; their separate exponents do not assert an
additional endpoint. This combination creates no new order of Appl, Caus or
orientation and no general licence for causativising every relational state.

Constructed continuous \ird{svornéra} \trsl{makes bigger} and
retrospective \ird{svornéranav} \trsl{has made bigger} follow the same
NOM/ERG case contrast. Their forms preserve the same expanded causative and the ordinary
conditioning of outer Voice.

Stative applicatives require a lexically licensed relation to the state.
The stimulus construction with \ird{všihn-} \trsl{be happy} relates an
experiencer to the news about which the state holds. Its applied stimulus is
direct and its experiencer transitive in the perfective:

\ex
\begingl
\gla Senica Marku všihnébik.//
\glb news.\Dir{} Marek.\Trs{} be.happy-\Appl{}-\Pf{}//
\glft \trsl{Marek was happy about the news on that occasion.}
\lingcontext{lexically restricted stimulus applicative}//
\endgl
\xe

The perfective bounds the state on the supplied occasion, while the applied
stimulus identifies what Marek is happy about. It does not itself introduce Become. Neither this stimulus relation nor
perfective aspect asserts that the news caused the state to arise; a causal
relation requires causative morphology or an independently established lexical
causative.
The stimulus extension supplies no general rule admitting a causer with every
stative applicative.

Underived stative verbs are intransitive and lack the lexical patient required
by the antipassive. Their derived causatives and applicatives license additional
participants through their respective argument-building heads; those
derivations do not establish a patient-bearing reading of the stative base.

Stative verbs may further function as attributive modifiers of nouns and remain
the productive attributive strategy in Iridian. Alongside them stands the small
closed set of paired adnominal property-root forms just mentioned; this use is
described in Chapter~\ref{ch:nouns}.

Property-root-based stative stems also feed the productive derived manner
adverbs in \ird{-ě} (§\,\ref{sec:stative-adv}), yielding lexical adverbial
adjuncts such as \ird{skorně} \trsl{quickly}. These belong to the clause's
preverbal adjunct structure, not to the attributive system proper.

\subsection{Syntactic analysis of stative verbs}
\label{sec:vm-tension-stative}

The lower stativising head projects \synhead{StateP}, whose holder or
experiencer is licensed by the predicate's lexical structure. Overt
\synhead{BecomeP} selects that state and introduces change, while restricted
\synhead{ManifestP} derives its outward display; both belong below inner
valency, Caus and orientation and require their separately listed input licences. The resulting predicate is selected by ordinary
Asp.Align. The schematic containment relation is
\[
\begin{aligned}
\mathrm{StateP}&\subset(\mathrm{BecomeP}/\mathrm{ManifestP})\subset
(\mathrm{ApplP}/\mathrm{VoiceP}_{\mathrm{AP}})\\
&\subset(\mathrm{CausP})\subset\mathrm{VoiceP}_{\mathrm{orient}}
\subset\mathrm{Asp.AlignP}.
\end{aligned}
\]
Parentheses mark independently licensed structure, not freely combinable
slots: in particular, underived states lack the patient frame required for
antipassive. The established actor-introducing layer and the restrictions on
inner Appl/Voice retain their positions in the general clause spine.

A lower \synhead{ResP} represents an event-derived state when the
construction independently entails a preceding event. Lexical
\ird{hvodna} \trsl{is open} entails openness without entailing an opening
event; a result-state interpretation of an independently licensed opening
predicate entails that event as well as its relevant state. ResP has no new
universal exponent here. Retrospective relevance alone does not force a
persisting result state, and a lexical state is not automatically reanalysed
as ResP. More specific target/resultant-state distinctions require separate
constructional evidence and are not assigned by terminology alone.

The distinction between a finite stative predicate and a nominal predicate is
also reflected in negation, even where neither clause contains an overt copula:

\begin{center}
\begin{tabular}{l l l}
\toprule
\textsc{predication type} & \textsc{affirmative} & \textsc{negative} \\
\midrule
Zero-copula equative & \ird{Marek doktor} & \ird{Marek doktor staly} \\
Stative verbal & \ird{Nužo mužna} & \ird{Nužo zamužna} \\
\bottomrule
\end{tabular}
\end{center}

\medskip
\noindent The nominal predicate \ird{doktor} takes clause-final
\ird{staly}, whereas \ird{mužna}, formed from \ird{mužn-} with continuous
\ird{-a}, takes \ird{za-}. Both clauses lack an overt copula, but the stative
predicate bears verbal and aspectual morphology.

Neg is analysed as selecting a nominal predication in the equative and an
aspectually inflected verbal predicate in the stative clause
(§\,\ref{sec:vm-neg}). Finite property predicates contain the verbalising
suffix \ird{-n-} and inflect for Asp.Align; predicative use of underived
property forms is undocumented.

The syntactic infrastructure of the stative domain therefore distinguishes
three outputs from the same lexical base. First, the finite stative verb heads
the clause's verbal predicate and requires no copula in ordinary property
predication. Second, a class-selected ordinary ATTR form is a non-finite derivative
inside the \synhead{DP}, corresponding to the adnominal strategy discussed in
Chapter~\ref{ch:nouns}; it is not a finite predicate and cannot bear aspect or
alignment. Third, the \ird{-ě} form is a lexical adverbial derivative that
functions as a clausal adjunct in the preverbal field, not as an inflectional
member of the verbal paradigm. The same lexical base thus feeds clause-level
predication, \synhead{DP}-internal attribution, and adverbial modification through three
distinct syntactic realisations.

Because underived stative predicates are intransitive, the alignment component
of Asp.Align is trivial in their simple finite use: the sole argument is the
direct-case absolutive/nominative argument in every aspect. The ordinary
alignment alternations re-enter the system only when valency-changing
morphology is added. Causativised statives behave like other derived
transitives, showing nominative alignment in the progressive and contemplative
and ergative alignment in the perfective and retrospective. The stative system
thus does not challenge the broader verbal theory; it instantiates it with a
lexically stative or overtly inchoative base whose participant structure is
built below the established case and pivot domain. A single-participant
inchoative remains direct in every aspect; Become neither inserts transitive
orientation nor bypasses the case machinery.


\section{Copular and nominal predication}\label{sec:copula}

Iridian distributes equative and classificatory predication across two
constructions that differ both structurally and semantically. The default
strategy is the \term{zero copula}: a null functional head that licenses a bare
equative predication without event structure. The second strategy is the overt
marked copular verb \ird{neh-}, which introduces an event of transition and
projects a full verbal spine. Whereas the zero copula expresses present or
atemporal classification, \ird{neh-} presents entry into that classification
as an event with aspectual structure.

Property predication --- the expression of adjectival meanings such as \trsl{be
beautiful} or \trsl{be good} --- belongs to the stative verb system
described in §\,\ref{sec:stative}. Existential
predication through the \ird{jec-}/\ird{zad-} pair is structurally distinct and
treated separately in §\,\ref{sec:existential-constructions}. The equative
strategies, their structural type, negation, and temporal range are summarised
in Table~\ref{tab:copula-overview}.

\begin{table}[ht]
  \footnotesize
  \caption{Equative and classificatory predication strategies.}\label{tab:copula-overview}
  \begin{tabularx}{\textwidth}{L L L L}
    \toprule
    \textsc{predication type} & \textsc{construction} &
    \textsc{negation} & \textsc{temporal range} \\
    \midrule
    Equative / class-membership & zero copula
      ($\varnothing$Cop\textsuperscript{0}) &
      \ird{staly} (free morpheme) & present / atemporal \\
    Transitional / inchoative & \ird{neh-} (overt verb) &
      \ird{za-neh-} (prefix) & full aspect paradigm \\
    Property / stative & stative verb &
      \ird{za-} + stative verb & full aspect paradigm \\
    \bottomrule
  \end{tabularx}
\end{table}

\section{Zero-copula predication}%
\label{sec:zero-copula}

\subsection{The construction and its distribution}

The default pattern for equative clauses --- those asserting identity, role, or
class membership --- is zero copula. The topic \synhead{DP} occupies the initial
position and the predicative complement appears in the direct case in
clause-final position; no overt verb is present. This is the stylistically
neutral, unmarked construction.

\pex
\a
\begingl
\gla Marek doktor.//
\glb Marek.\Dir{} doctor.\Dir{}//
\glft \trsl{Marek is a doctor.}//
\endgl
\a
\begingl
\gla Bách němscvení kompozitor.//
\glb Bach.\Dir{} German.\textsc{rel} composer.\Dir{}//
\glft \trsl{Bach is a German composer.}//
\endgl
\xe

Inclusive and class-membership predications follow the same zero-copula pattern
and carry no inference beyond simple predication. The predicative complement
always appears in the direct case regardless of which case would otherwise be
assigned in the same structural position.

\subsection{Structural analysis: \synhead{EquP} and the absence of AspP}

The zero copula is analysed as the phonologically null exponent of
Cop\textsuperscript{0}, which projects an \term{equative phrase}
(\synhead{EquP}) licensing the relation between the topic \synhead{DP} and
the predicative \synhead{DP}. In the unmarked present or atemporal
construction, \synhead{EquP} contains no aspectual projection and is selected
by \synhead{NegP}, when negation is present, and then by any selected
phasal or comment-domain layers before the higher \synhead{IP} layer. With
no such operators, \synhead{IP} selects the equative predication directly. Transitional interpretations instead employ \ird{neh-},
whose event structure supports the verbal aspect paradigm
(§\,\ref{sec:neh}); property predication likewise belongs to the stative
verbal system. The affirmative zero-copula structure is represented below:

\begin{center}
\begin{forest}
  syntax tree
  [\synhead{EquP}
    [{\synhead{DP}\\\textit{Marek}\\{\footnotesize\textsc{[dir]}}\\{\scriptsize topic}}]
    [{Eq$'$}
      [{\synhead{DP}\\\textit{doktor}\\{\footnotesize\textsc{[dir]}}\\{\scriptsize pred.}}]
      [{Eq$^{0}$\\{$\varnothing$}}]
    ]
  ]
\end{forest}
\end{center}

The two DPs bear the direct case as the default case assigned under
\synhead{EquP}: both the topic and the predicative complement are direct-case
nominals within a caseless, event-free predication domain.

\subsection{Temporal restriction}

The temporal corollary of the absence of \synhead{AspP} is that the zero-copula
construction is restricted to \term{present-time} and \term{atemporal} or
\term{generic} predications. A sentence like \ird{Marek doktor} is interpreted
as holding at the moment of utterance or as a classification that transcends any
particular time. It cannot, without additional context, be interpreted as a past
or future statement. This restriction follows directly from the lack of any
tense or aspect projection: there is no functional head in the clause that could
anchor the predication to a non-present time.

For past equative predications, \ird{neh-} in the retrospective (\ird{nav}) is
used; see §\,\ref{sec:neh-past}. For future-oriented predictions, the
contemplative of \ird{neh-} (\ird{náž}) is available. The zero copula's
restriction to present and generic contexts thus parallels the zero-copula
pattern found in Russian, Arabic, and Hebrew, where the copula is similarly
absent in the present indicative but required in other tenses --- though in
Iridian the structural basis is the absence of \synhead{EquP}'s eventuality
rather than the suppression of a tensed copula.

Negation of zero-copula predication uses clause-final \ird{staly}; its
structure and interaction with phasal particles are described in
§\,\ref{sec:negative-equatives}.

\section{Identificational equatives and participant nominals}%
\label{sec:ident-equative}\label{sec:nmlz-predicative}

When the predicative complement of a zero-copula equative clause is a
\term{participant nominal} (a nominalised verb form,
§\,\ref{sec:nmlz-predicative}), the construction is
\term{identificational}: the topic \synhead{DP} is identified with the referent
of the nominal description, which is interpreted as definite in the following example.
The English definite descriptions here and in the related existential examples
presume contextually recoverable referents; neither unanchored indirect
\ird{lúbcí} nor participant \ird{-ou} itself encodes definiteness:

\pex
\a
\begingl
\gla Lobek Janku pjaštník.//
\glb apple.\Dir{} Janek.\Trs{} eat-\textsc{trs}-\Pf{}//
\glft \trsl{Janek ate the apple.} \lingcontext{verbal clause: propositional assertion}//
\endgl
\a
\begingl
\gla Lobek Jancí pjaštníkou.//
\glb apple.\Dir{} Janek.\Ind{} eat-\textsc{trs}-\Pf{}-\Nmlz{}//
\glft \trsl{The apple is what Janek ate.} \lingcontext{identificational equative: patient nominal}//
\endgl
\xe

\noindent In both clauses, bare \ird{lobek} permits the contextually
recoverable reading expressed by English \trsl{the apple}; shared-anchored
\ird{lúbko} would explicitly mark retrieval through the common discourse
model (§\,\ref{sec:nm-def}). That anchoring contrast leaves the predication
type unchanged: transitive \ird{Janku} with finite \ird{pjaštník} forms the
verbal clause, whereas indirect \ird{Jancí} with participant nominal
\ird{pjaštníkou} forms the equative. In (b), the apple is identified with the
referent of the nominalised description. An
agent-nominal equative identifies the topic as the agent of the nominalised
event:

\ex
\begingl
\gla Janek lúbcí pjáščikou.//
\glb Janek.\Dir{} apple.\Ind{} eat-\Antip{}-\Pf{}-\Nmlz{}//
\glft \trsl{Janek is the one who ate the apple.}//
\endgl
\xe

\noindent The indirect dependents belong to the nominal case domain:
\ird{Jancí} specifies the agent of the patient description, while
\ird{lúbcí} specifies the theme associated with the antipassive-based agent
description. The latter is licensed by the nominal relational head described
in §\,\ref{sec:nominalised}; it is not an overt patient argument inside
antipassive \synhead{VP}. The same nominal derivation appears in existential
pivots, independently of matrix case or predication type.

Identificational negation rejects the proposed participant identification;
its interpretation and its contrast with finite verbal negation are illustrated
in §\,\ref{sec:negative-identificational-constructions}.

\section{The overt copular verb \ird{neh-}: transition and inchoative
predication}\label{sec:neh}

\subsection{Status and structural contrast with the zero copula}

The root \ird{neh-} is the overt copular verb of Iridian. Its most basic
function is to encode \term{transition}: the topic \synhead{DP} is asserted to
be moving from one classification or condition into another, or to have crossed
a threshold between them. This transitional meaning is not a pragmatic
implication but a semantic contribution of the root itself.

Zero-copula \synhead{EquP} lacks \synhead{AspP}, whereas the transition
introduced by \ird{neh-} can be viewed with respect to its onset, duration,
or completion. Accordingly, Asp\textsuperscript{0} selects the verbal
\synhead{VP}, and a negative clause projects \synhead{NegP} above that
aspectual predicate. The two structures are represented below:

\begin{center}
\begin{minipage}[t]{0.42\textwidth}
\centering
\textit{Zero copula:}\\[6pt]
\begin{forest}
  syntax tree
  [\synhead{EquP}
    [{\synhead{DP}\\{\footnotesize topic}}]
    [{Eq$'$}
      [{\synhead{DP}\\{\footnotesize pred.}}]
      [{Eq$^{0}$\\{$\varnothing$}}]
    ]
  ]
\end{forest}
\end{minipage}%
\hfill
\begin{minipage}[t]{0.52\textwidth}
\centering
\textit{Overt copula:}\\[6pt]
\begin{forest}
  syntax tree
  [\synhead{AspP}
    [{\synhead{DP}\\{\footnotesize topic}}]
    [{Asp$'$}
      [\synhead{VP}
        [{\synhead{DP}\\{\footnotesize pred.}}]
        [{\ird{neh-}}]
      ]
      [{Asp$^{0}$}]
    ]
  ]
\end{forest}
\end{minipage}
\end{center}

\noindent The consequence for negation is immediate: since \ird{neh-} is a
verbal root embedded under \synhead{AspP}, the standard verbal negation strategy
applies. Neg\textsuperscript{0} is spelled out as the prefix \ird{za-}, which
procliticizes to the left edge of the inflected form. \ird{Staly} does not arise
because the condition for its appearance --- the absence of a verbal host --- is
not satisfied.

\subsection{Paradigm}

The paradigm of \ird{neh-} is strongly suppletive and morphologically irregular;
principal forms are given in Table~\ref{tab:neh-paradigm}. Contemplative
\ird{náž} retains the copula's suppletive stem selection while carrying
the long contemplative exponent; its written final \ird{ž} has the regular
word-final devoiced realisation [ʃ].

\begin{table}[ht]
  \footnotesize
  \caption{Principal forms of \ird{neh-} (overt copula).}\label{tab:neh-paradigm}
  \begin{tabularx}{0.84\textwidth}{L l l}
    \toprule
    {\sc form} & {\sc affirmative} & {\sc negative} \\
    \midrule
    Continuous          & \ird{neha}   & \ird{zaneha} \\
    Progressive         & \ird{nejme}  & \ird{zanejme} \\
    Contemplative       & \ird{náž}   & \ird{zanáž} \\
    Perfective          & \ird{nek}    & \ird{zanek} \\
    Retrospective       & \ird{nav}   & \ird{zanav}\\
    \bottomrule
  \end{tabularx}
\end{table}

\noindent The negative forms are entirely regular: unstressed \ird{za-}
procliticizes onto each consonant-initial affirmative form without displacing
its primary stress. Before a vowel-initial host the allomorph is \ird{z-}
(§\,\ref{sec:neg-allomorphy}). No suppletive negative root exists for \ird{neh-}; the
formation parallels that of all other finite verbs.

\subsection{Aspectual semantics}

\ird{Neh-} introduces a transition event, and its aspectual forms characterise
different perspectives on that event in the usual way.

The \term{continuous} (\ird{neha}) presents the transition as ongoing and
unbounded: the subject is in the process of moving into the new classification.
The inchoative flavour is inherent to \ird{neh-} itself, not to the continuous
aspect.

\pex
\a
\begingl
\gla Marek doktor neha.//
\glb Marek.\Dir{} doctor.\Dir{} \Cop{}-\Cont{}//
\glft \trsl{Marek is becoming a doctor.}//
\endgl
\a
\begingl
\gla Marek doktor zaneha.//
\glb Marek.\Dir{} doctor.\Dir{} \Neg{}-\Cop{}-\Cont{}//
\glft \trsl{Marek is not becoming a doctor.}//
\endgl
\xe

The \term{progressive} (\ird{nejme}) adds an agentive or effortful dimension:
the subject is actively engaged in the process of transition. Where the
continuous reports a change in progress, the progressive foregrounds the subject's
directed activity toward the new state.

\ex
\begingl
\gla Marek doktor nejme.//
\glb Marek.\Dir{} doctor.\Dir{} \Cop{}-\Prog{}//
\glft \trsl{Marek is working toward becoming a doctor.}//
\endgl
\xe

The \Ctpl{} (\ird{náž}) locates the transition in prospective
time: the subject will enter the new classification. It is neutral between
imminent and more remote futures.

\ex
\begingl
\gla Marek doktor náž.//
\glb Marek.\Dir{} doctor.\Dir{} \Cop{}-\Ctpl{}//
\glft \trsl{Marek will become a doctor.}//
\endgl
\xe

The \term{perfective} (\ird{nek}) presents the transition as a completed
event: the subject has crossed the threshold and entered the new state. The
resulting state is typically understood as still holding unless contextually
cancelled.

\pex
\a
\begingl
\gla Tak byl kapitan nek.//
\glb \Dem{}.\Prox{} child.\Dir{} captain.\Dir{} \Cop{}-\Pf{}//
\glft \trsl{This child became captain.}//
\endgl
\a
\begingl
\gla Tak byl kapitan zanek.//
\glb \Dem{}.\Prox{} child.\Dir{} captain.\Dir{} \Neg{}-\Cop{}-\Pf{}//
\glft \trsl{This child did not become captain.}//
\endgl
\xe

\subsection{The retrospective and past equative predication}%
\label{sec:neh-past}

The \term{retrospective} (\ird{nav}) presents the transition as situated in a
past frame and viewed from the current moment. Its semantics are aspectually
neutral between the transitional reading (the subject entered the state at some
point in the past) and the stative reading (the subject was in the state during
the relevant past period).

\pex
\a
\begingl
\gla Marek doktor nav.//
\glb Marek.\Dir{} doctor.\Dir{} \Cop{}-\Ret{}//
\glft \trsl{Marek was a doctor. / Marek had been a doctor.}//
\endgl
\a
\begingl
\gla Marek doktor zanav.//
\glb Marek.\Dir{} doctor.\Dir{} \Neg{}-\Cop{}-\Ret{}//
\glft \trsl{Marek was not a doctor.}//
\endgl
\xe

This dual function of the retrospective is a direct consequence of the temporal
restriction on the zero copula. Since the zero copula is confined to
present-time and atemporal predications, past equative contexts require an overt
verb, and \ird{nav} is the form that fills this role. The transitional
implication that \ird{neh-} normally carries is therefore weakened or absent in
retrospective equatives: \ird{Marek doktor nav} does not necessarily imply that
Marek underwent a process of becoming a doctor, only that the doctor
classification held in the past. The implication is a pragmatic default, not a
semantic entailment of the retrospective form, and is readily cancelled by
context.

Whether the past classification still holds is not encoded. \ird{Marek doktor
nav} is compatible both with Marek's remaining a doctor and with his having
ceased to be one; the current-state inference, if present, is a conversational
implicature from the choice of a past-framed form over the zero-copula present.

The sequential converb \ird{nesu} --- colloquially \ird{nu} --- is a suppletive
form established independently of the regular aspect paradigm. Its
grammaticalisation as a discourse particle is described in
Section~\ref{sec:seq-copula} of Chapter~\ref{ch:complex}.

\section{Existential constructions}
\label{sec:existential-constructions}

\subsection{Overview}
\label{sec:existential-overview}%
\label{sec:jec}% retained for cross-reference compatibility
\label{sec:zad}% retained for cross-reference compatibility
\label{sec:jec-zad}% retained for cross-reference compatibility

Iridian expresses existential predication through a pair of suppletive verbal
roots that together constitute the affirmative/negative alternation of a single
functional head, Exist\textsuperscript{0}. The affirmative root \ird{jec-}
asserts the existence or presence of a referent; the negative root \ird{zad-}
denies it. Both inflect for the full aspect system; both are morphologically
irregular and suppletive, with the perfective and retrospective built on
distinct suppletive stems. The two paradigms are given together in
Table~\ref{tab:jec-zad-paradigm}. The contemplatives select
\ird{jec-} and \ird{zad-} with long \ird{-áž}; their irregularity in
other aspectual cells does not require a separate contemplative ending.

\begin{table}[ht]
  \footnotesize
  \caption{Principal forms of \ird{jec-} (\textsc{exst}) and \ird{zad-} (\textsc{neg.exst}).}\label{tab:jec-zad-paradigm}
  \begin{tabularx}{0.88\textwidth}{L l l}
    \toprule
    {\sc form} & \ird{jec-} (\textsc{exst}) & \ird{zad-} (\textsc{neg.exst}) \\
    \midrule
    Continuous          & \ird{jeca}   & \ird{zada} \\
    Progressive         & \ird{jecime} & \ird{zdíme} \\
    Contemplative       & \ird{jecáž} & \ird{zadáž} \\
    Perfective          & \ird{jak}    & \ird{zděk} \\
    Retrospective       & \ird{ján}   & \ird{znení} \\
    \bottomrule
  \end{tabularx}
\end{table}

The affirmative and negative roots are analysed as suppletive allomorphs of
Exist\textsuperscript{0}, with \ird{jec-} in affirmative clauses and
\ird{zad-} in the environment of Neg. Neg\textsuperscript{0} receives a zero
exponent while remaining present as Neg\textsuperscript{0}$[u\textsc{neg}]$ in the syntax;
the silent operator in Spec,\synhead{NegP} contributes semantic negation
(§\,\ref{sec:neg-exst-alternation}; §\,\ref{sec:vm-jec-zad}).

Iridian has three main productive uses of the existential construction: to
assert existence or presence, including location in a particular place
(\S\,\ref{sec:existence-and-presence}), to express possession
(\S\,\ref{sec:existence-possession}), and certain idiomatic presentational
constructions (\S\,\ref{sec:existential-presentatives}).


\subsection{Existence and presence}\label{sec:existence-and-presence}

The primary use of \ird{jec-} is to assert that a referent exists or is present
at the moment of utterance or within a contextually given frame. The entity in
question stands in the direct case. \ird{Zad-} forms the direct negative
counterpart, asserting non-existence or absence under the same structural
conditions.

\pex
\a
\begingl
\gla Byl jeca.//
\glb child.\Dir{} \Exst{}-\Cont{}//
\glft \trsl{There is a child. / A child is present.}//
\endgl
\a
\begingl
\gla Byl zada.//
\glb child.\Dir{} \NegExst{}-\Cont{}//
\glft \trsl{There is no child. / No child is present.}//
\endgl
\xe

\pex
\a
\begingl
\gla Kroumaště na pivo so jeca.//
\glb refrigerator.\Ind{} in beer.\Dir{} still \Exst{}-\Cont{}//
\glft \trsl{There is still (some) beer left in the refrigerator.}//
\endgl
\a
\begingl
\gla Kroumaště na pivo zada.//
\glb refrigerator.\Ind{} in beer.\Dir{} \NegExst{}-\Cont{}//
\glft \trsl{There is no beer in the refrigerator.}//
\endgl
\xe

Within the Minimalist framework, \ird{jec-} is the phonological exponent of a
dedicated functional head, \term{Exist}\textsuperscript{0}, which projects
\term{ExistP} as the innermost stratum of the clause spine --- the slot occupied
by root \synhead{VP} in ordinary verbal clauses. Its sole argument is the \term{existential pivot}, the direct-case
\synhead{DP} whose existence or presence is asserted. ExistP is then
selected by Asp\textsuperscript{0}, which contributes the aspect suffix, and
feeds the higher functional sequence including \synhead{NegP}. When
\synhead{NegP} is absent, Exist\textsuperscript{0} is spelled out as \ird{jec-};
when \synhead{NegP} is present, suppletive Vocabulary Insertion replaces
\ird{jec-} with \ird{zad-}, leaving Neg\textsuperscript{0} with no independent
PF exponent. The following diagrams show only the existential/aspectual constituent and,
in the negative, its negation layer. Finite inflection and optional higher
operators and ordinary placement are omitted; the illustrated pivot
positions are lower chain occurrences when the full finite placement field
is projected. The diagrams do not place existential Asp above those operators. The affirmative and negative structures of \ird{Byl jeca} and
\ird{Byl zada} are shown below:

\begin{center}
\begin{forest}
  syntax tree
  [\synhead{AspP}
    [{\synhead{DP}\\\textit{byl}\\{\footnotesize child.\Dir{}}\\{\scriptsize pivot}}]
    [{Asp$'$}
      [ExistP
        [{Exist$'$}
          [{Exist\\\ird{jec-}}]
        ]
      ]
      [{Asp\\\ird{-a}\\{\footnotesize\textsc{[cont]}}}]
    ]
  ]
\end{forest}
\hspace{2em}
\begin{forest}
  syntax tree
  [\synhead{NegP}
    [{$\mathrm{Op}_{\neg}$\\{\footnotesize$[i\textsc{neg}]$}}]
    [{Neg$'$}
      [\synhead{AspP}
        [{\synhead{DP}\\\textit{byl}\\{\footnotesize child.\Dir{}}\\{\scriptsize pivot}}]
        [{Asp$'$}
          [ExistP
            [{Exist$'$}
              [{Exist\\\ird{zad-}\\{\scriptsize suppletive}}]
            ]
          ]
          [{Asp\\\ird{-a}\\{\footnotesize\textsc{[cont]}}}]
        ]
      ]
      [{Neg\\$\varnothing$\\{\footnotesize$[u\textsc{neg}]$}}]
    ]
  ]
\end{forest}
\end{center}

\noindent The direct-case pivot precedes the predicate in
Spec,\synhead{AspP} in both structures; in the negative, the silent operator
occupies Spec,\synhead{NegP} above it. This low pivot position does not itself
confer aboutness topichood: in a
presentational context, the pivot may introduce a new discourse referent
without occupying Spec,\synhead{TopP} (§\,\ref{sec:tc-topicless}).
Negative clauses contain a zero exponent
of Neg\textsuperscript{0}, which conditions the suppletive root \ird{zad-}
at Exist\textsuperscript{0}.

\noindent\textit{Location.}\enspace \ird{Jec-} is also used to locate an entity: the entity
appears in the direct case and the location is expressed by an indirect noun
phrase followed by a postposed clitic adposition. \ird{Zad-} negates location,
asserting that the referent is not present at the named location.

\pex
\a
\begingl
\gla Marek dumě na jeca.//
\glb Marek.\Dir{} house.\Ind{} in \Exst{}-\Cont{}//
\glft \trsl{Marek is at home.}//
\endgl
\a
\begingl
\gla Marek dumě na zada.//
\glb Marek.\Dir{} house.\Ind{} in \NegExst{}-\Cont{}//
\glft \trsl{Marek is not at home. / Marek is absent.}//
\endgl
\xe

\noindent The locative phrase is merged as a PP adjunct to ExistP, introducing
the spatial domain within which the existential relation holds, without altering
the valency of Exist\textsuperscript{0}. The pivot remains the sole direct
argument; the structures parallel those for bare existentials above, with the PP
adjoined at the ExistP periphery in both affirmative and negative derivations.


\section{Possessive predication}\label{sec:existence-possession}

Iridian has no separate verb meaning \trsl{have}. What English presents as
\trsl{I have X} is in Iridian construed as \trsl{there is X at me}: the
possessor stands in the indirect case, the possessed thing in the direct, and
the construction is a special case of existential predication. \ird{Zad-}
produces the negative possessive, asserting that the possessed item does not
exist within the possessor's sphere.

The ordinary arrangement of the clause places the possessor phrase first, the
possessed noun next, and the existential verb last:

\begin{center}
[\emph{possessor}.\Ind{} \ \ird{na}] \ \ [\emph{possessee}.\Dir{}] \ \ \ird{jec-} / \ird{zad-}
\end{center}

\noindent The possessed noun is the thing whose existence is asserted or denied
relative to the possessor as reference point. In grammatical terms it is the
existential pivot --- the direct argument of Exist\textsuperscript{0} --- while
the possessor is an oblique locative domain, a PP adjunct to ExistP, within
which the existential relation holds.

\pex
\a
\begingl
\gla Bylam na tóm jeca.//
\glb child.\Ind{} in/on book.\Dir{} \Exst{}-\Cont{}//
\glft \trsl{The child has a book.}//
\endgl
\a
\begingl
\gla Bylam na tóm zada.//
\glb child.\Ind{} in/on book.\Dir{} \NegExst{}-\Cont{}//
\glft \trsl{The child does not have a book.}//
\endgl
\xe

\pex
\a
\begingl
\gla Tóm bylam na jeca.//
\glb book.\Dir{} child.\Ind{} in/on \Exst{}-\Cont{}//
\glft \trsl{As for the book, it is with the child.}//
\endgl
\a
\begingl
\gla Tóm bylam na zada.//
\glb book.\Dir{} child.\Ind{} in/on \NegExst{}-\Cont{}//
\glft \trsl{As for the book, the child does not have it.}//
\endgl
\xe

\noindent The second pair illustrates that topic-initial order is fully
available with \ird{zad-} on the same terms as with \ird{jec-}: the alternation
belongs to the preverbal field and is independent of the affirmative/negative
opposition.

Following Freeze (1992) and the localist tradition, possession is treated here
as a structural subtype of location: \trsl{the book exists at the child}. The
possessor phrase [\textit{possessor}.\Ind{} \ird{na}] and a plain locative
phrase [\textit{location}.\Ind{} \ird{na}] are structurally parallel: both are
oblique PPs that modify ExistP without altering the valency of
Exist\textsuperscript{0}. The affirmative and negative possessive structures
are:

\begin{center}
\begin{forest}
  syntax tree
  [ExistP
    [{PP\\\ird{bylam na}\\{\footnotesize possessor}}]
    [ExistP
      [{\synhead{DP}\\\textit{tóm}\\{\footnotesize book.\Dir{}}\\{\scriptsize pivot}}]
      [{Exist$'$}
        [{Exist\\\ird{jec-}}]
      ]
    ]
  ]
\end{forest}
\hspace{2em}
\begin{forest}
  syntax tree
  [\synhead{NegP}
    [{$\mathrm{Op}_{\neg}$\\{\footnotesize$[i\textsc{neg}]$}}]
    [{Neg$'$}
      [\synhead{AspP}
        [{Asp$'$}
          [ExistP
            [{PP\\\ird{bylam na}\\{\footnotesize possessor}}]
            [ExistP
              [{\synhead{DP}\\\textit{tóm}\\{\footnotesize book.\Dir{}}\\{\scriptsize pivot}}]
              [{Exist$'$}
                [{Exist\\\ird{zad-}\\{\scriptsize suppletive}}]
              ]
            ]
          ]
          [{Asp\\\ird{-a}\\{\footnotesize\textsc{[cont]}}}]
        ]
      ]
      [{Neg\\$\varnothing$\\{\footnotesize$[u\textsc{neg}]$}}]
    ]
  ]
\end{forest}
\end{center}

\noindent The structural parallelism between locative and possessive
existentials is a direct consequence of this unified analysis: both
constructions merge a PP at the ExistP periphery, and the only grammatical
difference is the semantic content of the oblique phrase itself. The suppletive
readjustment applies identically in both frames; the PP adjunct is unaffected by
the presence or absence of \synhead{NegP}.

Since possession is expressed through the existential verb, possessive clauses
naturally inflect across the full aspectual system, and \ird{jec-} and
\ird{zad-} participate symmetrically in these aspectual contrasts. The
continuous is the common unmarked form: \ird{Bylam na tóm jeca} \trsl{the child
has a book}; \ird{Bylam na tóm zada} \trsl{the child does not have a book}. The
progressive conveys a relation in the course of arising: \ird{Bylam na tóm
jecime} \trsl{the child is getting a book}. The contemplative projects into
prospect: \ird{Bylam na tóm jecáž} \trsl{the child will have a book}. The
perfective of this lexically alternating existential permits an acquisition
reading:
\ird{Bylam na tóm jak} \trsl{the child got a book}. The retrospective places the
relation in a prior frame: \ird{Bylam na tóm žav} \trsl{the child had a book}.
In all these forms the construction remains fundamentally existential: the
possessed noun remains the direct pivot, the possessor remains oblique, and the
aspect construes the lexically available relation or acquisition event.
This existential alternation does not make perfective aspect a general
inchoative operator on productive property statives
(§\,\ref{sec:stative-inchoative}).

Since the possessed noun is the existential pivot, indefinite possessees are
especially natural. Definite possessees are not excluded, but with them the
clause more readily shades toward a meaning of location, availability, or
custody rather than settled ownership --- a consequence of the existential
source of the construction.


\section{Event-anchored existential pivots}\label{sec:existential-presentatives}

The existential predicate in Iridian consistently selects a nominal pivot: the
entity asserted to exist stands as the direct-case nominal argument of
Exist\textsuperscript{0}. The nominal pivots admitted by this selectional requirement fall into
three subtypes according to the internal structure of the pivot \synhead{DP}:

\begin{enumerate}[wide]
\item a bare \synhead{NP} (plain existential; §\,\ref{sec:existence-and-presence}),
\item a complex \synhead{DP} with an overt lexical head modified by a
  participant-oriented attributive (overt-head presentational;
  §\,\ref{sec:exst-overt-head}),
\item a headless \synhead{DP} whose sole content is a participant nominal
  derived from verbal structure by Iridian's participant-nominal morphology
  (headless participant-nominal; §\,\ref{sec:exst-headless-nmlz}).
\end{enumerate}

\noindent The three subtypes share a single structural identity as the nominal
argument of Exist\textsuperscript{0}, originating in Spec,\synhead{ExistP}
and appearing in the low Spec,\synhead{AspP} in the bare and presentational
aspectual structures illustrated here; the PP-initial possessive retains the
pivot in the lower existential domain. They differ in the internal architecture of the pivot
\synhead{DP}; neither low position confers aboutness topichood, and both fall
within the scope of $\mathrm{Op}_{\neg}$ in a negative existential.
Ordinary finite pronunciation precedes the particle field through the
placement dependency; PP-initial possessives retain their independently
specified PP--pivot order rather than acquiring a new aboutness topic.

\subsection{The participant-nominal pivot}
\label{sec:exst-pivot-arch}

An unidentified event participant may be introduced by a participant nominal,
which combines nominal external distribution with a verbal event description
(§\,\ref{sec:nmlz-predicative}). In the analysis
adopted here, $n_{\mathrm{part}}$ supplies nominal category and participant
abstraction directly, without an intervening null lexical noun. Negative
indefinites such as \ird{zide} \trsl{no one} and
\ird{deno} \trsl{nothing} provide a distinct lexical strategy under strict
negative concord (§\,\ref{sec:nm-negative-concord}); their availability does not
require a silent noun within participant-nominal pivots.

The pivot's orientation follows from the same embedded aspect and Voice
configuration as in other participant nominals
(§\,\ref{sec:nm-derived-participant}). With transitive bases, perfective
and retrospective agents require antipassive Voice, whereas active continuous,
progressive, and contemplative forms already privilege the agent. An
intransitive nominal selects its sole participant without antipassive
morphology. Patient-oriented Voice, realised as \ird{-n-} or as
\ird{-an-} after a rhotic causative, and applicative promotion retain their
ordinary selectional roles, subject to the higher causative configuration where present. Existential
predication introduces no additional operation selecting an event participant.

\subsection{Overt-head presentational pivots}
\label{sec:exst-overt-head}

The overt-head presentational pivot consists of a lexical head noun in the direct
case modified by a participant-oriented attributive. Its canonical shape is:

\begin{center}
[\textit{(nominal dependent.\Ind{})} \ \textit{participant-oriented-relative} \ \textit{head-NP}.\Dir{}] \ \ \ird{jeca}
\end{center}

The modifier selects the participant privileged by embedded Voice and aspect.
Non-agentive heads are preferred in this presentational use, so patient-oriented
and applicative relatives are especially natural, as in \trsl{a letter that
the woman wrote} or \trsl{a child for whose benefit someone ate}. This preference
concerns the choice of pivot description; it does not change the relative's participant
selection.

\ird{Jeca} is fixed in this use. Temporal and aspectual interpretation come
entirely from the embedded attributive, which may be perfective, retrospective,
continuous, progressive, or contemplative. The matrix \ird{jeca} signals only
that the referent is being introduced into the current discourse.

\pex
\a
\begingl
\gla Kářnikový senica jeca.//
\glb read-\textsc{trs}-\Pf{}-\textsc{rel} news.\Dir{} \Exst{}-\Cont{}//
\glft \trsl{There is some news that (I) read.} \lingcontext{patient/theme introduced}//
\endgl
\a
\begingl
\gla Že dbórnikový kíl jeca.//
\glb \textsc{1sg}.\Ind{} hear-\textsc{trs}-\Pf{}-\textsc{rel} song.\Dir{} \Exst{}-\Cont{}//
\glft \trsl{There is a song that I heard.} \lingcontext{stimulus introduced via perception event}//
\endgl
\a
\begingl
\gla Védnikový ryc jeca.//
\glb see-\textsc{trs}-\Pf{}-\textsc{rel} cat.\Dir{} \Exst{}-\Cont{}//
\glft \trsl{There is a cat that (I) saw.} \lingcontext{perceived referent introduced via completed event}//
\endgl
\a
\begingl
\gla Pjaštébikový byl jeca.//
\glb eat-\Appl{}-\Pf{}-\textsc{rel} child.\Dir{} \Exst{}-\Cont{}//
\glft \trsl{There is a child for whose benefit someone ate.} \lingcontext{promoted participant introduced}//
\endgl
\xe

The attributive is an adjunct within the pivot \synhead{DP}; it is not a
predicate of the matrix clause and contributes no independent propositional
content at the matrix level. Exist\textsuperscript{0} asserts the existence
of an individual satisfying the complex pivot description. The construction
fixes the matrix form as \ird{jeca}, while the embedded relative supplies
its own event viewpoint; this restriction does not extend to the locative and
possessive existentials, which admit aspectual variation.

\begin{center}
\begin{forest}
  syntax tree
  [\synhead{AspP}
    [{\synhead{DP}\\{\footnotesize pivot}}
      [{RelP\\\ird{kářnikový}\\{\footnotesize read-\textsc{trs}-\Pf{}-\textsc{rel}}}]
      [{\synhead{NP}\\\ird{senica}\\{\footnotesize news.\Dir{}}}]
    ]
    [{Asp$'$}
      [ExistP
        [{Exist$'$}
          [{Exist\\\ird{jec-}}]
        ]
      ]
      [{Asp\\\ird{jec-a}\\{\scriptsize Exist+Cont}}]
    ]
  ]
\end{forest}
\end{center}

Patients, themes, recipients, and other non-agentive heads are strongly
favoured in the overt-head presentational, whereas an agent-oriented relative
is marked in this frame:

\ex
\begingl
\gla ?Seničete kážikový chrás jeca.//
\glb news.\Ind{}.\Def{} read-\textsc{antip}-\Pf{}-\textsc{rel} person.\Dir{} \Exst{}-\Cont{}//
\glft \trsl{There is a person who read the news.} \lingcontext{formally possible, but
  marked: agentive head disfavoured in overt-head presentational}//
\endgl
\xe

This discourse preference is, however, a property of the overt-head subtype
specifically. The headless participant-nominal construction
(§\,\ref{sec:exst-headless-nmlz}) is the default indefinite strategy across all
participant roles and does not inherit any preference for non-agentive pivots.

The negative counterpart of the overt-head presentational is formed with
\ird{zada}. \synhead{NegP} scopes over ExistP, and the construction asserts that
no referent satisfying the event-anchored description exists. As with
\ird{jeca}, the continuous form \ird{zada} is fixed; the aspectual
interpretation of the embedded event comes entirely from the attributive.

\pex
\a
\begingl
\gla Kářnikový senica zada.//
\glb read-\textsc{trs}-\Pf{}-\textsc{rel} news.\Dir{} \NegExst{}-\Cont{}//
\glft \trsl{There is no news that (I) read.} \lingcontext{negative overt-head presentational; event-anchored referent denied}//
\endgl
\a
\begingl
\gla Pjaštébikový byl zada.//
\glb eat-\Appl{}-\Pf{}-\textsc{rel} child.\Dir{} \NegExst{}-\Cont{}//
\glft \trsl{There is no child for whose benefit someone ate.} \lingcontext{negative presentational; promoted participant denied}//
\endgl
\xe

\begin{center}
\begin{forest}
  syntax tree
  [\synhead{NegP}
    [{$\mathrm{Op}_{\neg}$\\{\footnotesize$[i\textsc{neg}]$}}]
    [{Neg$'$}
      [\synhead{AspP}
        [{\synhead{DP}\\{\footnotesize pivot}}
          [{RelP\\\ird{kářnikový}\\{\footnotesize read-\textsc{trs}-\Pf{}-\textsc{rel}}}]
          [{\synhead{NP}\\\ird{senica}\\{\footnotesize news.\Dir{}}}]
        ]
        [{Asp$'$}
          [ExistP
            [{Exist$'$}
              [{Exist\\\ird{zad-}\\{\scriptsize suppletive}}]
            ]
          ]
          [{Asp\\\ird{zad-a}\\{\scriptsize NegExst+Cont}}]
        ]
      ]
      [{Neg\\$\varnothing$\\{\footnotesize$[u\textsc{neg}]$}}]
    ]
  ]
\end{forest}
\end{center}

\noindent The negative existential clause denies that any referent satisfies the
event-anchored description expressed by the complex pivot. The
attributive remains an adjunct within the pivot \synhead{DP} with no
matrix-level propositional force. The disfavour of agentive pivot heads carries
over to \ird{zada} for the same reasons as in the affirmative.

The \ird{zad-} existential takes a singular pivot; plural participant nominals
are not licensed in the negative existential frame.

\subsection{Headless participant-nominal pivots}
\label{sec:exst-headless-nmlz}

The headless participant-nominal construction employs a pivot \synhead{DP} that
contains no lexical head noun. The pivot's entire content is a participant
nominal derived from verbal structure, whose voice--aspect configuration determines
participant orientation. This construction is the default means of expressing
indefinite participant reference in Iridian.

The participant nominal has the same $n_{\mathrm{part}}$--Asp--Voice
structure as the identificational predicates of §\,\ref{sec:nmlz-predicative}, with construction-conditioned exponence and an
independent nominal case domain. The participant is selected by the complete
voice--aspect configuration (Table~\ref{tab:participant-nominals}); the
nominaliser neither assumes finite absolutive assignment nor promotes an
otherwise unprivileged argument.

\noindent When such a nominal occupies the existential pivot position,
Exist\textsuperscript{0} asserts the existence of an individual satisfying the
participant description. The event description is internal to the nominalised
pivot. With contemplative aspect, the existence of a participant does not
assert realisation of the projected event. The matrix \ird{jec-}
contributes only $\textsc{exist}(x)$; the verbal material, together with its
aspect and voice, supplies the property that identifies the individual
introduced.

\pex
\a
\begingl
\gla Ščénikou jeca.//
\glb arrive-\Pf{}-\Nmlz{} \Exst{}-\Cont{}//
\glft \trsl{There is someone who arrived.} \lingcontext{sole-participant indefinite: intransitive nominal}//
\endgl
\a
\begingl
\gla Lúbcí pjáščikou jeca.//
\glb apple.\Ind{} eat-\Antip{}-\Pf{}-\Nmlz{} \Exst{}-\Cont{}//
\glft \trsl{There is someone who ate the apple.} \lingcontext{agentive indefinite: antipassive agent-nominal}//
\endgl
\a
\begingl
\gla Védnikou jeca.//
\glb see-\textsc{trs}-\Pf{}-\Nmlz{} \Exst{}-\Cont{}//
\glft \trsl{There is something (I) saw.} \lingcontext{patient indefinite: transitive patient-nominal}//
\endgl
\a
\begingl
\gla Pjaštébikou jeca.//
\glb eat-\Appl{}-\Pf{}-\Nmlz{} \Exst{}-\Cont{}//
\glft \trsl{There is someone for whose benefit someone ate.} \lingcontext{promoted participant: applicative nominal}//
\endgl
\xe

The formal structure of the headless participant-nominal existential differs
from the overt-head construction only in the internal architecture of the pivot.
The pivot contains \synhead{AspP} above \synhead{VoiceP}, with the root
and its verbal categoriser in the lowest constituent.
The head $n_{\mathrm{part}}$ introduces nominal category and participant
abstraction; no lexical noun is present. Optional causative structure lies
below the outer patient-selecting Voice head and above the lower predicate,
which may itself contain applicative or antipassive morphology
(§\,\ref{sec:morph-template}). A simple applicative privileges its applied
participant; a higher causative instead determines the selected participant
according to its derived frame. Existential predication contributes no
additional event-participant selection.

\begin{center}
\begin{forest}
  syntax tree
  [\synhead{AspP}
    [{\synhead{DP}\\{\footnotesize pivot}}
      [{\synhead{nP}}
        [\synhead{AspP}
          [\synhead{VoiceP}
            [{\synhead{vP}\\{Root + categoriser}}]
            [{Voice}]
          ]
          [{Asp}]
        ]
        [{n\textsuperscript{0}\\\ird{-ou}\\{\scriptsize\Nmlz{}}}]
      ]
    ]
    [{Asp$'$}
      [ExistP
        [{Exist$'$}
          [{Exist\\\ird{jec-}}]
        ]
      ]
      [{Asp\\\ird{jec-a}\\{\scriptsize Exist+Cont}}]
    ]
  ]
\end{forest}
\end{center}

\noindent The semantics is parallel to the plain and overt-head cases. The nP
denotes a property of individuals, event-anchored by the verbal material it
encodes; Exist\textsuperscript{0} introduces a witness individual satisfying
that property. The eventive content is entirely internal to the pivot; no
clausal embedding of the verbal structure into the matrix occurs.

The contrast between the two constructional subtypes is already visible in the
grammar's existential question--answer paradigm:

\pex
\a
\begingl
\gla Védnikou by jeca?//
\glb see-\textsc{trs}-\Pf{}-\Nmlz{} \textsc{q} \Exst{}-\Cont{}//
\glft \trsl{Did you see anything?} \lingcontext{headless NMLZ pivot; referent as yet unidentified}//
\endgl
\a
\begingl
\gla Védnikový ryc jeca.//
\glb see-\textsc{trs}-\Pf{}-\textsc{rel} cat.\Dir{} \Exst{}-\Cont{}//
\glft \trsl{There is a cat that I saw.} \lingcontext{overt-head pivot; specific referent introduced}//
\endgl
\xe

\noindent In (a), headless \ird{védnikou} supplies the participant property
whose satisfaction is questioned. In (b), \ird{ryc} additionally classifies
the participant as a cat, while \ird{védnikový} supplies the event-related
restriction. The question leaves the participant unidentified; the answer
introduces a referent under a lexical description. The presence of the lexical
noun does not itself encode specificity.

\subsection{Indefinite participant introduction}
\label{sec:exst-indef-pos}

The headless participant-nominal construction is the primary means by which
Iridian expresses what English achieves with the free indefinite pronouns
\trsl{somebody}, \trsl{someone}, \trsl{something}, and \trsl{anything}. The
existential pivot is a participant nominal whose voice--aspect configuration selects the
relevant participant role.

\pex
\a
\begingl
\gla Ščénikou jeca.//
\glb arrive-\Pf{}-\Nmlz{} \Exst{}-\Cont{}//
\glft \trsl{Someone arrived.} \lingcontext{lit.\ there is one-who-arrived; sole-participant indefinite}//
\endgl
\a
\begingl
\gla Lúbcí pjáščikou jeca.//
\glb apple.\Ind{} eat-\Antip{}-\Pf{}-\Nmlz{} \Exst{}-\Cont{}//
\glft \trsl{Someone ate the apple.} \lingcontext{lit.\ there is one-who-ate-it; agentive transitive indefinite}//
\endgl
\a
\begingl
\gla Védnikou jeca.//
\glb see-\textsc{trs}-\Pf{}-\Nmlz{} \Exst{}-\Cont{}//
\glft \trsl{(I) saw something.} \lingcontext{lit.\ there is a seen-thing; patient indefinite}//
\endgl
\a
\begingl
\gla Pjaštébikou jeca.//
\glb eat-\Appl{}-\Pf{}-\Nmlz{} \Exst{}-\Cont{}//
\glft \trsl{Someone benefited from an eating event.} \lingcontext{lit.\ there is one-for-whose-benefit-someone-ate; applied participant}//
\endgl
\xe

The headless nominals in (a) and (b) freely introduce the sole participant
of arriving and the actor of eating, respectively. The preference against
agent-oriented heads belongs to the overt-head presentational; it does not
restrict participant reference in the headless construction.

In these existential clauses, the participant nominal constitutes the pivot
without an unexpressed lexical head noun, while embedded Voice and aspect
specify its participant role.

\subsection{Negative existential indefinites}
\label{sec:exst-indef-neg}

Negative existential clauses with participant-nominal pivots express the
meanings rendered in English by \trsl{nobody}, \trsl{no one}, and
\trsl{nothing}. The derivation is structurally parallel to the affirmative,
with \synhead{NegP} added above the aspectual predicate. Its silent operator
$\mathrm{Op}_{\neg}$$[i\textsc{neg}]$ takes scope over the existential predication, while
Neg\textsuperscript{0}$[u\textsc{neg}]$ conditions suppletive Vocabulary Insertion at
Exist\textsuperscript{0}. The resulting clause denies that any individual
satisfies the nominal description. A participant nominal formed from
affirmative verbal material does not acquire a negative-concord feature merely
by serving as the pivot: its negative indefinite interpretation follows from
existential predication within the scope of negation.

\pex
\a
\begingl
\gla Ščénikou zada.//
\glb arrive-\Pf{}-\Nmlz{} \NegExst{}-\Cont{}//
\glft \trsl{Nobody arrived.} \lingcontext{lit.\ there is no one-who-arrived}//
\endgl
\a
\begingl
\gla Lúbcí pjáščikou zada.//
\glb apple.\Ind{} eat-\Antip{}-\Pf{}-\Nmlz{} \NegExst{}-\Cont{}//
\glft \trsl{Nobody ate the apple.} \lingcontext{lit.\ there is no one-who-ate-it}//
\endgl
\a
\begingl
\gla Védnikou zada.//
\glb see-\textsc{trs}-\Pf{}-\Nmlz{} \NegExst{}-\Cont{}//
\glft \trsl{Nothing was seen. / (I) saw nothing.} \lingcontext{lit.\ there is no seen-thing}//
\endgl
\xe

The negative existential denies that any individual satisfies the participant
description, without presupposing an established event. Identificational
negation, by comparison, rejects a participant identification against a
contextually given event (§\,\ref{sec:exst-vs-ident}).

Neither negative existential \ird{zada} nor identificational \ird{staly}
negates the embedded verbal description. Participant nominalisation selects
actual \synhead{AspP}, including inner valency, causative and orientation
structure, and hence excludes internal verbal \ird{za-} as well as higher
Phas, Comm, Mood and clausal Evid structure (§\,\ref{sec:nominalised}).
The matrix structure $[\mathrm{NegP}\ [\mathrm{AspP}_{\mathrm{matrix}}\ [\mathrm{ExistP}\
[n_{\mathrm{part}}\ [\mathrm{AspP}\ \ldots]]]]]$ thus differs from excluded
$*[n_{\mathrm{part}}\ [\mathrm{NegP}\ \ldots]]$. The same distinction
separates a matrix phasal/comment operator over an existential or equative
predication from an operator inside its participant description. A negative
existential with
\ird{pjáščikou} denies the existence of an eater satisfying the description;
it does not introduce someone characterised as not having eaten.

\subsection{Relation to identificational equatives}
\label{sec:exst-vs-ident}

The same participant nominals that serve as existential pivots with
\ird{jec-}/\ird{zad-} also serve as predicates in identificational equatives
(§\,\ref{sec:ident-equative}). As an existential pivot, the nominal supplies the description whose
satisfiers are asserted to exist; as an equative predicate, it supplies the
description with which the topic is identified.

In an identificational equative, the participant nominal is the predicative
complement of a zero-copula clause, and the topic is identified with its referent.
In an identificational reading with a contextually established event,
negation with \ird{staly} rejects the topic as the relevant participant
while retaining the event as background:

\pex
\a
\begingl
\gla Janek lúbcí pjáščikou.//
\glb Janek.\Dir{} apple.\Ind{} eat-\Antip{}-\Pf{}-\Nmlz{}//
\glft \trsl{Janek is the one who ate the apple.} \lingcontext{identificational equative: event presupposed; Janek identified as agent}//
\endgl
\a
\begingl
\gla Janek lúbcí pjáščikou staly.//
\glb Janek.\Dir{} apple.\Ind{} eat-\Antip{}-\Pf{}-\Nmlz{} \Neg{}//
\glft \trsl{Janek is not the one who ate the apple.} \lingcontext{identificational negation: event still presupposed; Janek misidentified as agent}//
\endgl
\xe

\noindent The suggestion that someone else ate the apple depends on the
continued assumption that there was an eater, together with the rejection of
Janek as that person (§\,\ref{sec:ident-equative}).

With the same nominal as an existential pivot, the affirmative asserts the
existence of an eater, while the negative denies it:

\pex
\a
\begingl
\gla Lúbcí pjáščikou jeca.//
\glb apple.\Ind{} eat-\Antip{}-\Pf{}-\Nmlz{} \Exst{}-\Cont{}//
\glft \trsl{Someone ate the apple.} \lingcontext{existential: asserts existence of an agent; no event presupposition}//
\endgl
\a
\begingl
\gla Lúbcí pjáščikou zada.//
\glb apple.\Ind{} eat-\Antip{}-\Pf{}-\Nmlz{} \NegExst{}-\Cont{}//
\glft \trsl{Nobody ate the apple.} \lingcontext{negative existential: no agent exists; event not presupposed}//
\endgl
\xe

\noindent In the existential construction, the participant nominal supplies the
pivot, and the negative clause denies that any individual satisfies its
description, without requiring the event to be established in advance.
The existential assertion thus differs from an identification whose event and
participant description are already given in the discourse.

To assert that no participant satisfies a description, the existential
construction with \ird{zad-} supplies the relevant reading; the identificational
construction with \ird{staly} instead disputes the proposed identity of a
participant. In contexts where the event is already established, the latter
naturally preserves that background and suggests a different participant,
whereas the negative existential makes no such commitment.

Lexical negative indefinites such as \ird{deno} \trsl{nothing} and
\ird{zide} \trsl{no one} participate in strict negative concord with both
existential and equative negation (§\,\ref{sec:nm-negative-concord}). Their $[u\textsc{neg}]$ feature requires the local negative
operator, whose presence is expressed by \ird{zad-} or \ird{staly} according
to predicate type; neither the indefinite nor the overt negative exponent
introduces another semantic negation. This lexical strategy coexists with the
participant-nominal existential, whose pivot need not contain a concord item.

\subsection{Discourse profile of the two constructional subtypes}
\label{sec:exst-discourse}

The overt-head presentational introduces a referent under both a lexical
category and an event-related description. The headless construction supplies
the participant description alone, making it suitable when the participant is
otherwise unidentified. Overt-head pivots favour non-agentive referents,
whereas headless pivots show no corresponding restriction by participant role.

The two constructions may therefore form a question--answer sequence:
\ird{Védnikou by jeca?} \trsl{Did you see anything?} asks whether an
individual satisfies the participant description, and \ird{Védnikový ryc
jeca} \trsl{There is a cat that I saw} introduces such an individual as a
cat. The same answer is poorly suited to \trsl{Who saw the cat?}, which asks
for the agent of an already specified event rather than for the patient
introduced by this pivot. The contrast follows from the participant selected
by the relative and the information requested in the question; it does not
make agents generally incompatible with existential predication.

\section{Comparative and superlative constructions}\label{sec:comparison}

Comparison in Iridian is expressed analytically using particles that combine
with stative verbs. The subject of comparison stands in the direct case; the
standard of comparison takes the indirect case. Two particles are
distinguished: \ird{stá} for direct comparison (\textit{more\ldots than},
\textit{-er than}) and \ird{zrc} for equative comparison (\textit{as\ldots
as}).

\subsection{Direct comparison}

Direct comparison is formed with the particle \ird{stá} placed between the
standard of comparison and the stative verb. The stative verb inflects for the
ordinary aspectual system within the comparative frame. A simple state stays
a state; an overt \ird{-él-}/\ird{-el-} transition can be viewed as
unfolding or bounded. Comparison specifies the scalar standard separately
from this state/change contrast:

\pex
\a
\begingl
\gla Janek Marcí stá četna.//
\glb Janek.\Dir{} Marek.\Ind{} \textsc{comp} be.tall-\Cont{}//
\glft \trsl{Janek is taller than Marek.}//
\endgl
\a
\begingl
\gla Janek Marcí stá četnélime.//
\glb Janek.\Dir{} Marek.\Ind{} \textsc{comp} be.tall-\textsc{inch}-\Prog{}//
\glft \trsl{Janek is becoming taller than Marek.}//
\endgl
\a
\begingl
\gla Janek Marcí stá četnélek.//
\glb Janek.\Dir{} Marek.\Ind{} \textsc{comp} be.tall-\textsc{inch}-\Pf{}//
\glft \trsl{Janek became taller than Marek.}//
\endgl
\xe

\noindent where \ird{četn-} is \trsl{be tall}. The particle \ird{stá} may also
combine with the \ird{-á} attributive form (§\ref{sec:stative-attr}) to form
an attributive comparative. The particle precedes the \ird{-á} form
prenominally:

\pex
\a \ird{stá četná pres} \trsl{the taller man}
\a \ird{Marcí stá četná pres} \trsl{the man who is taller than Marek}
\xe

\subsection{Equative comparison}

Equative comparison is formed with the particle \ird{zrc}. The structure is
otherwise identical to the direct comparative, with the subject in the direct
case and the standard in the indirect. Negation of the equative falls on the
stative verb itself, which takes the regular negative prefix \ird{za-}
(§\ref{sec:negation}):

\pex
\a
\begingl
\gla Janek Marcí zrc četna.//
\glb Janek.\Dir{} Marek.\Ind{} \textsc{eq} be.tall-\Cont{}//
\glft \trsl{Janek is as tall as Marek.}//
\endgl
\a
\begingl
\gla Janek Marcí zrc začetna.//
\glb Janek.\Dir{} Marek.\Ind{} \textsc{eq} \Neg{}-be.tall-\Cont{}//
\glft \trsl{Janek is not as tall as Marek.}//
\endgl
\xe
